% ===================================================================
%  TABELO N° 11 — La Urbo e lua monumenti. --- L'Incendio.
%  Feuillets 77 a 81 du fac-simile = folios 75 a 79.
%  Correspondance : folio imprime = numero de feuillet - 2.
%
%  Ca tabelo esas en du parti : « La Urbo e lua monumenti » (folios
%  75-78) e « L'Incendio » (folios 78-79), qua riprenas sua propra
%  numerado d'alineas de 1 — videz §6 dil consigne pri la clefi %%K.
% ===================================================================

% --- Feuillet 77 (folio 75) -----------------------------------------
\begin{VUpage}[77]{75}
%%K t11-apar-1 apar
\VUsaut{2.0mm}
\VUcentre{13.2pt}{90}{TRIESMA SERIO}
\VUsaut{3.0mm}
\VUfilet{16mm}
\VUsaut{5.6mm}
\VUtitre{15.0pt}{60}{TABELO N\textsuperscript{o} 11}
\VUsaut{1.4mm}
\VUpk{11.4pt}{40}{La Urbo e lua monumenti. --- L'Incendio.}
\VUsaut{2.2mm}

%%K t11-tit-1 sub
\VUpk{10.2pt}{40}{La Preurbo.}
\VUsaut{3.0mm}

%%K t11-c1-01-1 p
\VUblancAlinea
1. --- Me habitas granda urbo situita ye\nl
100 kilometri de Oceano, e qua tamen esas\nl
\VUgras{portuo} \textsuperscript{(1)} altaranga. La fluvio qua bordizas olu\nl
havas larjeso de 400 metri e formacas tre bela\nl
rado.

%%K t11-c1-02-1 p
\VUblancAlinea
2. --- La tota parto dil urbo qua esas sur la\nl
\VUgras{kayi} \textsuperscript{(2)} tote aspektas marala ed industriala, e\nl
formacas \VUgras{preurbo} \textsuperscript{(3)} qua esas habitata nur da\nl
navani e laboristi. Ici laboras en la \VUgras{fabrike}\cc
\VUgras{rii} \textsuperscript{(4)} ed en la \VUgras{gasiferio} \textsuperscript{(5)} di qua la alta \VUgras{kame}\cc
\VUgras{no} \textsuperscript{(6)} e la \VUgras{gasometro} videsas de tre fore.

%%K t11-c1-03-1 p
\VUblancAlinea
3. --- En la mezepoko, la urbo esis fortifikita,\nl
ed on ankore trovas kelka restaji di lua rem\cc
pari, aparte la \VUgras{pordo} \textsuperscript{(8)} dil olda \VUgras{fortifikita kas}\cc
\VUgras{telo} \textsuperscript{(7)} flankumita per lua du \VUgras{turmi} \textsuperscript{(9)} qui nun\nl
uzesas kom \VUgras{karcero.}

%%K t11-c1-04-1 p
\VUblancAlinea
4. --- En ica tota \VUgras{quartero,} qua aspektas tre\nl
olda, nur un konstrukturo esas relate moder\cc
na : to esas la \VUgras{doganeyo} \textsuperscript{(10)}. Ol esas situita in\cc
ter la kayo e la \VUgras{stradeto} \textsuperscript{(11)} qua duktas a la\nl
olda \VUgras{komon-domo} \textsuperscript{(12)}. On facile rikonocas ta
\parplein
\end{VUpage}

% --- Feuillet 78 (folio 76) -----------------------------------------
\begin{VUpage}[78]{76}
%%K t11-c1-04-1 p suite
\VUcontinue
lasta monumento pro la \VUgras{klosheyo} \textsuperscript{(13)} gigantala\nl
qua dominacas l'antiqua urbo.

%%K t11-c1-05-1 p
\VUblancAlinea
5. --- Sur l'altra latero di la strado, \VUgras{stacas}\nl
edifico konstruktita samstile kam la doganeyo:\nl
to esas la \VUgras{Borso} \textsuperscript{(14)}. Un de lua fasadi regardas\nl
mikra placo ube esas la \VUgras{prefekteyo} \textsuperscript{(15)}. Fakte,\nl
nia urbo sen esar chef-urbego, esas tamen im\cc
portanta chef-urbo di departmento.

%%K t11-tit-2 sub
\VUsaut{5.0mm}
\VUpk{10.2pt}{40}{La Maxim Alta Parto dil Urbo.}
\VUsaut{3.4mm}

%%K t11-c1-06-1 p
\VUblancAlinea
6. --- La garnizono, sat grandanombra, esas\nl
lojata en vasta \VUgras{kazerni} \textsuperscript{(16)} tre salubra; li ex\cc
tensas su til la gretaro dil \VUgras{parko komonala} \textsuperscript{(17)}.\nl
La \VUgras{bulvardo} \textsuperscript{(19)} qua duktas ad ica parko pasas\nl
dop la \VUgras{spar-kaseyo} \textsuperscript{(18)} ed abutas a la placo\nl
dil \VUgras{katedralo} \textsuperscript{(20)}.

%%K t11-c1-07-1 p
\VUblancAlinea
7. --- Ta omna monumenti su extensas am\cc
fiteatre sur kolineto ube esas anke nia vasta\nl
\VUgras{liceo} \textsuperscript{(21)}.

%%K t11-c1-07-2 p
\VUblancAlinea
Se on decensas tra voyo kelke inklinita, qua\nl
rajuntas la Granda strado, on arivas al\nl
\VUgras{Tribunalo} \textsuperscript{(22)} avan qua extensas su agreabla\nl
\VUgras{gardeno publika} \textsuperscript{(26)}. Ta \VUgras{garden-placo} ne\nl
esas tre granda, ma ol facas meze dil urbo\nl
oaziso de verdajo e koldeteso. Pro to ol\nl
esas tre frequentata da la urbani qui pri\cc
zas venar sideskar sub la tendo dil \VUgras{kafee}\cc
\VUgras{rio} \textsuperscript{(25)}, por askoltar la soldata muzikisti qui\nl
jovdie e sundie pleas en la \VUgras{kiosko} \textsuperscript{(27)} pozita\nl
inter la gazoneyi e masivi.

%%K t11-c1-08-1 p
\VUblancAlinea
8. --- Sur un de ta gazoneyi stacas bronza\nl
\VUgras{statuo} \textsuperscript{(28)}, monumento erektita por la memoro\nl
di un de nia samurbani, qua facis granda servi\nl
a sua naskal urbo.
\end{VUpage}

% --- Feuillet 79 (folio 77) -----------------------------------------
\begin{VUpage}[79]{77}
%%K t11-tit-3 sub
\VUsaut{4.0mm}
\VUpk{10.2pt}{40}{La Transportili.}
\VUsaut{3.4mm}

%%K t11-c1-09-1 p
\VUblancAlinea
9. --- Til nun nia urbo ne ja esas lumizata\nl
per la lumo elektrala, lu havas nur \VUgras{gas-bru}\cc
\VUgras{lili} \textsuperscript{(29)}. Ma la \VUgras{hika jurnalaro} kampanias por\nl
ke ta sistemo di lumizo esez plubonigata, \VUgras{quale}\nl
on ja perfektigis \VUgras{la sistemo di traciono di la}\nl
tramvoyi. \VUgras{La} olda veturi, \VUgras{tirata} da \VUgras{kavali,} esis\nl
praktike remplasata per \VUgras{elektral tramvoyi} \textsuperscript{(31)}\nl
qui rapide transportas la voyajanti de un ex\cc
tremajo dil \VUgras{urbo al altra,} po tre \VUgras{chipa} preco.

%%K t11-c1-10-1 p
\VUblancAlinea
10. --- Proxim la Judicieyo esas la \VUgras{Banko} \textsuperscript{(32)}\nl
ube on diskontas la komercal valori. La \VUgras{ban}\cc
\VUgras{ko-komisionisti} \textsuperscript{(33)} iras inkasar domicile la\nl
debaji.

%%K t11-tit-4 sub
\VUsaut{4.0mm}
\VUpk{10.2pt}{40}{Arto ed Instrukto.}
\VUsaut{3.4mm}

%%K t11-c1-11-1 p
\VUblancAlinea
11. --- La bel edifico qua stacas proxime,\nl
esas la \VUgras{muzeo.} Ol kune kontenas la \VUgras{bibliote}\cc
\VUgras{ko} \textsuperscript{(34)} dil urbo, la bela \VUgras{kolektaji pri natur}\cc
\VUgras{cienco} \textsuperscript{(35)} (Naturmuzeo) e gracioza \VUgras{piktur-mu}\cc
\VUgras{zeo} \textsuperscript{(36)} qua konstitucas splendid \VUgras{expozeyo} per\cc
mananta.

%%K t11-c1-11-2 p
\VUblancAlinea
Ol apertesas por la publiko dufoye singla-\nl
semane, ma la stranjeri darfas vizitar lu irga\cc
die po gratifiketo al \VUgras{gardisto} \textsuperscript{(37)}. La \VUgras{piktis}\cc
\VUgras{ti} \textsuperscript{(38)} venas ibe por laborar. On vidas li \VUgras{avan}\nl
lia tresto, kun \VUgras{planketo} en manuo, kopiante\nl
la famoza pikturi.

%%K t11-c1-12-1 p
\VUblancAlinea
12. --- Sur la altajo, dop la muzeo, on vides\cc
kas la \VUgras{kupolo} \textsuperscript{(40)} di la \VUgras{hospitalo} \textsuperscript{(39)} e la kons\cc
trukturi di la \VUgras{normal skolo} \textsuperscript{(41)} en qua esis\nl
exercata l'instruktisti di nia skoli komonala.
\end{VUpage}

% --- Feuillet 80 (folio 78) -----------------------------------------
\begin{VUpage}[80]{78}
\VUblancAlinea
Sur la placo dil Teatro esas \VUgras{posto-kontoro} \textsuperscript{(43)}\nl
instalita en \VUgras{domo privata} \textsuperscript{(42)}.

%%K t11-tit-5 sub
\VUsaut{5.0mm}
\VUpk{11.4pt}{40}{L'Incendio.}
\VUsaut{3.0mm}

%%K t11-tit-6 sub
\VUpk{10.2pt}{40}{Klamo pri Incendio.}
\VUsaut{3.4mm}

%%K t11-c2-01-1 p
\VUblancAlinea
1. --- Teroriganta rumoro difuzesas tra la\nl
urbo : incendio jus eventis en la teatro ! Ye\nl
l'instanto kande me arivas sur la \VUgras{placo} \textsuperscript{(96)}, ja\nl
la turbo esas asemblita sur la \VUgras{trotuaro} \textsuperscript{(46)} e\nl
sur la \VUgras{choseo} \textsuperscript{(47)}. La \VUgras{postisti} \textsuperscript{(44)} e la \VUgras{letris}\cc
\VUgras{ti} \textsuperscript{(45)} ekiras la posto-kontoro. \VUgras{Tramduktisto} \textsuperscript{(48)}\nl
devis haltigar sua (tram)-veturo por lasar pasar\nl
\VUgras{ambulanco-veturo} \textsuperscript{(50)} urbala, qua arivas kun\nl
\VUgras{kirurgiisto} \textsuperscript{(52)} e \VUgras{flegisto} \textsuperscript{(51)} por flegar \VUgras{vundi}\cc
\VUgras{to} \textsuperscript{(53)} adportita per \VUgras{kamilo} \textsuperscript{(54)} proxim la \VUgras{jur}\cc
\VUgras{nal-kiosko} \textsuperscript{(55)}.

%%K t11-c2-02-1 p
\VUblancAlinea
2. --- Infantria kompanio, qua esis retroi\cc
ranta de l'exerci, haltis por sekurigar l'ordinal\nl
servado kun la \VUgras{kavalkanta jendarmi urbala} \textsuperscript{(61)}.\nl
\VUgras{La kavalkisti,} di qui la kavali kabras, plu\nl
bone kam la \VUgras{soldati} \textsuperscript{(57)} o la \VUgras{jendarmi} \textsuperscript{(63)}, re\cc
troirigas la \VUgras{kuriozi} \textsuperscript{(56)}. Cetere la jendarmi esas\nl
tre okupata. Un de li indikas a la \VUgras{policisti} \textsuperscript{(64)}\nl
adube on devas transportar altra \VUgras{acidentito} \textsuperscript{(65)},\nl
pumpisto kurajoza falinta de skalo.

%%K t11-c2-03-1 p
\VUblancAlinea
3. --- \VUgras{L'oficiro} \textsuperscript{(66)} precizigas la loko ube on\nl
devas lokizar la \VUgras{vapor-pumpilo} \textsuperscript{(67)} arivanta.\nl
Vartante ta potenta mashino, il baterie instali\cc
gis \VUgras{manu-pumpilo} \textsuperscript{(68)} di qua la longa \VUgras{tubo} \textsuperscript{(69)}\nl
lansas aquo torentatre sur la fairo.

\VUblancAlinea
Sis pumpisti manovras ol, dum ke lia kama\cc
rado, qua tenas la \VUgras{spricilo} \textsuperscript{(70)}, direktas la \VUgras{spri}\cc
\VUgras{co} \textsuperscript{(71)} ad la centro dil incendio.
\end{VUpage}

% --- Feuillet 81 (folio 79) -----------------------------------------
\begin{VUpage}[81]{79}
%%K t11-c2-04-1 p
\VUblancAlinea
4. --- An l'altra latero dil teatro, on apogis\nl
la granda \VUgras{skalo} \textsuperscript{(72)}. Ol posibligas a la pumpisti\nl
proximeskar ad la flami qui spricas inter vor\cc
tici de nigra \VUgras{fumuro} \textsuperscript{(75)}.

%%K t11-tit-7 sub
\VUsaut{5.0mm}
\VUpk{10.2pt}{40}{Brava Prodaji.}
\VUsaut{3.4mm}

%%K t11-c2-05-1 p
\VUblancAlinea
5. --- \VUgras{L'incendio} \textsuperscript{(74)} naskis en la foyero dil\nl
\VUgras{teatro} \textsuperscript{(73)}, dum ke on repetis \VUgras{l'opero} di qua\nl
l'unesma \VUgras{reprezento} eventabus morge. Fortu\cc
noze nur poka \VUgras{spektanti} \textsuperscript{(76)} esis en la spekto\cc
salono. La kompatindi refujis sur la \VUgras{balko}\cc
\VUgras{no} \textsuperscript{(77)} ube kurajoza \VUgras{pumpisti} \textsuperscript{(86)} venas sokur\cc
sar li per skalo e kordegi.

%%K t11-c2-06-1 p
\VUblancAlinea
6. --- Sub la \VUgras{peristilo} \textsuperscript{(78)}, ube la \VUgras{afisho} \textsuperscript{(79)}\nl
informas pri la reprezento, on vidas la \VUgras{muzi}\cc
\VUgras{kisti} \textsuperscript{(81)} di \VUgras{l'orkestro} \textsuperscript{(80)} fugar, kunportante\nl
lia instrumenti : \VUgras{violino} \textsuperscript{(81)}, \VUgras{fluto} \textsuperscript{(82)}, \VUgras{violon}\cc
\VUgras{celo} \textsuperscript{(83)}, \VUgras{trombono} \textsuperscript{(84)}, \VUgras{tamburo} \textsuperscript{(85)}, e. c.. On\nl
esforcas salvar parto dil \VUgras{moblaro} \textsuperscript{(87)}.

%%K t11-c2-07-1 p
\VUblancAlinea
7. --- \VUgras{L'aktoruli} \textsuperscript{(89)} e \VUgras{l'aktorini} \textsuperscript{(88)}, \VUgras{kantis}\cc
\VUgras{tuli} e \VUgras{kantistini} mustis fugar kun sua teatral\nl
\VUgras{kostumo} \textsuperscript{(90)}. La tenoro explikas a \VUgras{jurnalisto} \textsuperscript{(92)}\nl
quale to eventis. \VUgras{La raportisto} quik adkuris de\nl
l'unesma instanto kun la autoritatozi : \VUgras{pre}\cc
\VUgras{fekto} \textsuperscript{(91)}, \VUgras{urbestro} \textsuperscript{(93)}, \VUgras{komisario di polico} \textsuperscript{(94)}\nl
e \VUgras{judicial oficisto} \textsuperscript{(95)}, qui inquestos por judi\cc
ciar pri la responsivesi.

\VUsaut{6.0mm}
\VUfilet{22mm}
\end{VUpage}
